\section{Capacity stability of the graded current and energy calculus}
\label{sec:r11-graded-capacity-stability}

Use the domain, background and power-energy notation of
Sections~\ref{sec:r11-power-all-degree} and
\ref{sec:r11-first-order-power-currents}.
The limits below are ordinary original-current limits,
including every compact smooth complementary test and
every actual moment refinement. No assertion about
additional GK delta tests is implicit.

\begin{lemma}[Capacity-controlled coefficient multipliers]
\label{lem:r11-general-coefficient-capacity-multiplier}
On a fixed ordinary compact test stage, suppose
coefficient-Radon currents $P_l$ converge to $P$ on
original smooth tests. Suppose that for each fixed
complementary test $\eta$, all its coefficient variations,
including that of $P$, are bounded by $C_\eta c(A)$ on
Borel sets $A$, where $c$ is the same specified local
relative capacity. Let $a_l,a$ be actual bounded Borel
weights, with one common bound, such that $a_l\to a$ in
that local capacity and $a$ is quasicontinuous for its
ordinary-open outer version. Then
\[
                         a_lP_l\longrightarrow aP
 \tag{GS1}
\]
on the original compact tests of that stage.
Positivity and closedness of the $P_l$ are not required.
\end{lemma}

\noindent\textbf{Status: \texttt{DIRECT\_COROLLARY}} of
the frozen capacity multiplier argument, stated for
coefficient variation rather than only positive products.

\begin{proof}
For a fixed $\eta$, split the changing-weight error over
$\{|a_l-a|>\delta\}$ and its complement. Its variation
is at most
$2M C_\eta c(\{|a_l-a|>\delta\})$ plus $\delta$ times
the uniformly bounded compact coefficient mass.
Next remove an ordinary open set $G$ of arbitrarily small
capacity on whose complement $a$ is continuous.
On the compact remainder extend this bounded continuous
function to a bounded continuous function on the compact
stage. The latter extension is available in the original
metrizable affinoid neighborhood. Both errors over $G$
are bounded by a fixed multiple of $c(G)$, for the
approximating and limiting coefficient measures separately.
Original smooth scalar functions are uniformly dense in
continuous functions on the compact stage, by the frozen
original separation and cutoff argument. Their test
convergence and the uniform coefficient mass bound give
convergence with this continuous extension. Let the
capacity error and then $\delta$ tend to zero.
The assumed bound for the limiting current has been used
explicitly; it is not inferred by restricting a weak
limit to an arbitrary Borel set.
\end{proof}

\begin{lemma}[Bounded weighted first-order and energy stability]
\label{lem:r11-bounded-first-energy-capacity-stability}
Let bounded classical tuples converge in the exact
localized capacity sense of (PE16), with a common actual
bound. Over the fixed arbitrary original background,
their bounded mixed currents, first-order currents
$d'(v_iQ_\alpha)$ and $d''(v_iQ_\alpha)$, and energy
currents $\Gamma_{Q_\alpha}(v_i,v_a)$ converge on all
original compact tests. Each conclusion remains valid
after multiplication by a continuous function of the
bounded tuple.
\end{lemma}

\noindent\textbf{Status: \texttt{NEW\_PROVED}} for the
first-order and energy extension; the bounded mixed
current assertion is the frozen input already identified
in Lemma~\ref{lem:r11-bounded-analytic-capacity-stability}.

\begin{proof}
The bounded mixed stability theorem first gives
$Q_{\alpha,l}\to Q_\alpha$ and all continuous tuple
weights needed here. For clarity, a continuous function
on the common compact range box is uniformly continuous,
so its changing-weight gap is contained in a finite
union of potential gap sets. It is quasicontinuous off
a union of the finitely many bounded-potential exceptions.
Lemma~\ref{lem:r11-general-coefficient-capacity-multiplier}
therefore proves its weighted convergence, using the
bounded mixed capacity domination, without restricting
the scalar composition to one variable.

In particular $v_{i,l}Q_{\alpha,l}\to v_iQ_\alpha$.
Differentiating distributionally proves unweighted
first-order convergence. Estimate (FO4), applied to
the common bound on each local stage, controls the
coefficient variations of these derivatives and their
limit by the required local capacity.
Apply the preceding multiplier lemma to prove the
weighted first-order assertion.

For energies use the identity
\[
 \Gamma_{Q_\alpha}(v_i,v_a)
 =\tfrac12D(v_iv_aQ_\alpha)
       -\tfrac12v_iD(v_aQ_\alpha)
       -\tfrac12v_aD(v_iQ_\alpha).
 \tag{GS2}
\]
The weighted bounded mixed theorem passes all three
terms in distributions, hence gives unweighted current
convergence of the energies. Estimate (PG8), and
polarization of the positive energy matrix, bound
their coefficient variations by the same local
capacity. The multiplier lemma proves the last
weighted assertion. Derivatives of nonsmooth weights
have not been used anywhere in this proof.
\end{proof}

\begin{theorem}[Nonmonotone graded stability]
\label{thm:r11-nonmonotone-graded-stability}
Let $v_{i,l}\in\mathcal E^\sigma_{S,\partial}(V)$,
$1\leq i\leq m$, have a uniform bound on their
$\sigma$-power energies. Let $u_i$ be actual negative
extended classical psh functions satisfying
$\mathrm{BG}(u_i,0)$, and suppose
$v_{i,l}\to u_i$ in (PE16) in every slot.
Then all $u_i$ belong to the same power class.

Let $F$ be continuous on the nonpositive orthant with
$|F(x)|\leq A(1-\sum_i x_i)^a$, $a\geq0$.
With $k=|\alpha|$, the following convergences hold
under their respective strict inequalities:
\[
\begin{array}{ll}
 F(\mathbf v_l)N_\alpha(\mathbf v_l)
       \longrightarrow F(\mathbf u)N_\alpha(\mathbf u),
       &a+k<q+\sigma,\quad k\leq q,\\[2pt]
 F(\mathbf v_l){\cal J}'_{\alpha,i}(\mathbf v_l)
       \longrightarrow F(\mathbf u){\cal J}'_{\alpha,i}(\mathbf u),
       &a+k+1<q+\sigma,\quad k\leq q-1,\\[2pt]
 F(\mathbf v_l)G^\alpha_{ih}(\mathbf v_l)
       \longrightarrow F(\mathbf u)G^\alpha_{ih}(\mathbf u),
       &a+k+2<q+\sigma,\quad k\leq q-1.
\end{array}
 \tag{GS3}
\]
There is the corresponding double-prime statement.
Only terms whose indices meet the printed inequalities
are asserted.

Equality in any of the three exponent conditions is
also permitted if the following additional uniform
tail condition holds for $w_l=\sum_i v_{i,l}$:
\[
 \lim_{R\to\infty}\sup_l
       \int_R^\infty t^{q+\sigma-1}
              C_S(\{w_l<-t\};V)\,dt=0.
 \tag{GS4}
\]
This endpoint condition is not inferred from a
uniform bound on the full power energies alone.
\end{theorem}

\noindent\textbf{Status: \texttt{NEW\_PROVED}}.
The topology is the exact all-stage capacity topology
(PE16). Under the standing one-affinoid condition (DG1),
Corollary~\ref{cor:r11-affinoid-local-capacity-topology}
now proves its equivalence with fixed-domain capacity
convergence on compact cores. This is the explicit analytic
tester-extension theorem, not an assumed domain identification.

\begin{proof}
Apply Theorem~\ref{thm:r11-power-capacity-stability} to
each slot, repeated in the $q$ differentiated positions,
to obtain membership of the limits before constructing
their unbounded currents. The energy cone estimate
(PE10) bounds $e_\sigma(w_l)$ uniformly; (PE4) therefore
bounds the full integrals $ {\cal J}_{w_l}(0)$ uniformly.

For each fixed truncation level $j$, the preceding
bounded stability lemma applies to every original
compact test, including the bounded tuple weight.
Equations (PG3), (FO8), and (PG12) bound the errors
from removing the truncations, respectively, by
constants times
\[
 j^{a+k-q-\sigma},\qquad
 j^{a+k+1-q-\sigma},\qquad
 j^{a+k+2-q-\sigma},
 \tag{GS5}
\]
uniformly in $l$. Their additional factor
$ {\cal J}_{w_l}(j/4)$ may be bounded by its full
integral. The finitely many compact core controls
are uniform by bounded CLN at level one.
The corresponding limiting errors tend to zero
by the fixed-tuple estimates already proved.
First pass $l\to\infty$ at fixed $j$, then let
$j\to\infty$. This proves (GS3) whenever the
displayed exponent is negative.

At exponent zero retain the tail factor instead
of bounding it by the full integral. Assumption
(GS4) then supplies the same two-limit argument.
This proof does not pass a moving Borel restriction
through weak convergence or presume equality of
raw and outer capacities.
\end{proof}

\begin{corollary}[Full endpoint continuity for decreasing power families]
\label{cor:r11-decreasing-graded-endpoints}
If $v_{i,l}\downarrow u_i\in
\mathcal E^\sigma_{S,\partial}(V)$ are actual negative
extended classical psh functions, then (GS3) holds
with weak inequalities at every endpoint.
In particular, for $\sigma=1$, the top energy
currents $G^\alpha_{ih}$ with $|\alpha|=q-1$
are jointly decreasing-continuous on all original tests.
\end{corollary}

\noindent\textbf{Status: \texttt{DIRECT\_COROLLARY}}
of the proved graded estimates and bounded local
decreasing capacity convergence.

\begin{proof}
Upward stability and (PE9) put all rows in the
same power class with uniformly bounded energies.
The analytic common-collar theorem gives (PE16)
for each fixed bounded truncation. Since
$w_l\geq w=\sum_i u_i$, their capacity tails
are bounded pointwise by the tail of $w$.
Thus (GS4) follows from the finite integral
$J_\sigma(w)$. Apply the endpoint assertion.
\end{proof}

\paragraph{Why an endpoint needs a hypothesis.}
The bounded shell family in
Proposition~\ref{prop:r11-power-energy-continuity-failure}
has uniformly unit $\sigma$-power energy and converges
in (PE16) to zero, while its critical weighted top
measures converge to a nonzero mass at the rigid
origin. This refutes removal of the strict exponent
or uniform-tail qualification in the top weighted
case of (GS3). It does not contradict decreasing
endpoint continuity with an actual power-class limit.
